\section{Related Work}
\label{sec:related}

% TODO (~1–1.5 pages)
%
% Structure:
%
% 1. Cultural heritage KGs and LOD initiatives
%    - Europeana LOD (EDM; large scale, but limited SPARQL endpoint and no graph browser)
%    - Smithsonian LOD, British Museum CIDOC-CRM
%    - ArCo: Italian CH KG (CIDOC-CRM based, interlinked with DBpedia/GVP) [Carriero et al. ISWC 2019]
%    - ICH-KG: Chinese intangible CH (automated population from ICH database)
%    - ArtGraph: WikiArt + DBpedia (artistic KG with graph-based discovery) [Castellano et al. 2022]
%
% 2. Multimodal CH knowledge bases — ArtKB [Blanco et al. ESWC 2026] is the closest contemporary work:
%    - CACAO ontology (CIDOC-CRM extension), built from Wikidata, 8M triples, 307K artefacts
%    - Combines RDF graph (GraphDB) + vector DB (Qdrant) + object storage (MinIO) + API gateway
%    - Use cases: Text2SPARQL, neural image retrieval, AI metadata suggestion, Text2RDF
%    - Key difference from GeMeA: small curated art collection (307K) vs. national cross-domain
%      corpus (65M); API-only vs. full web UI; no named graph provenance; no self-hosting docs
%
% 3. RDA/FRBR-based linked bibliographic data
%    - BIBFRAME (Library of Congress): RDA-aligned, bibliographic focus
%    - Linked bibliographic data initiatives (OCLC, DNB)
%    - mocho as a bridge between EDM and RDA (this work)
%
% 4. KG browsers and SPARQL interfaces
%    - Wikidata Query Service (public SPARQL; no CH-specific browsing)
%    - DBpedia Spotlight / LodLive (generic LOD browsers)
%    - DDB native UI: search only, no SPARQL, no graph navigation, no open data export
%
% 5. Comparison table (rows: GeMeA, ArtKB, Europeana LOD, ArCo, DDB native UI)
%    Columns: scale, source, ontology, SPARQL endpoint, web UI, multimodal, open/self-hostable
%
% 6. ORKG comparison: TBD

\todo{Related Work.}
